\documentclass[10pt,a4paper]{amsart}
\usepackage[margin=19mm]{geometry}
\usepackage{amsmath,amssymb,mathtools}
\usepackage[scr=rsfs,cal=euler]{mathalfa}
\usepackage{xfrac}
\usepackage[unicode,hidelinks,bookmarks=false]{hyperref}
\newcommand{\ie}{i.e.\ }
\newcommand{\cf}{cf.\ }
\newcommand{\resp}{respectively\ }
\newcommand{\Subsection}{\subsection}
\providecommand{\nobreakdash}{\nobreak\hbox{-}}
\setlength{\emergencystretch}{2em}
\multlinegap0pt
\usepackage{bm}
\usepackage{bbm}
\usepackage{dsfont}
\usepackage{xspace}
\usepackage{cancel}
\usepackage{mathtools}
\usepackage{amsthm}
\makeatletter
\@ifundefined{theoremN}{\newtheorem{theoremN}{Theorem}}{}
\makeatother
\title[A Geometric Approach to the Links-Quivers Correspondence II:]{A Geometric Approach to the Links-Quivers Correspondence II: Rational Links}
\author{Jonathan A. Higgins}
\date{Source: arXiv:2603.01312v1 (CC BY 4.0)}
\begin{document}
\maketitle

\noindent\emph{Source and licence.} A Geometric Approach to the Links-Quivers Correspondence II: Rational Links. Version arXiv:2603.01312v1 (CC BY 4.0), distributed under the Creative Commons Attribution 4.0 International licence, as recorded in the arXiv licence metadata for this exact version and shown on the abstract page https://arxiv.org/abs/2603.01312v1. Full rights notice: licenses/arxiv-2603.01312-CC-BY-4.0.txt.\par\smallskip

\noindent\textit{Editorial scope.} Counted unit: the Theoremn whose source label is \texttt{None} in the article (source0.tex lines 3488--3494, source label \texttt{None}), delivered together with the article's own complete original proof of it (source0.tex lines 3496--3544, one \texttt{proof} environment of 49 lines). No mathematics is written, repaired or re-derived: statement and proof are the article's own text line for line. Required context retained uncounted: none. Every definition, display and label of those lines is retained unchanged. Omitted from this package and not redistributed: 1--3487, 3495--3495, 3545--3549. Nothing inside a retained excerpt is omitted or altered: every retained body is delivered exactly as the article's lines are written, so no omission receipt of any kind accompanies this package. Citations: the retained excerpts carry no citation command at all, so no bibliography entry of the article is transcribed and no reference is redistributed. Reference adaptation: none was needed and none was made; every reference inside the delivered unit resolves inside it, so no unresolved reference or citation is produced. Printed slips kept literal: no printed word, symbol, hypothesis or number of the counted statement or of its proof was altered. Layout and class: the article's own class is \texttt{amsart} with packages including graphicx, url, tabularx, array, geometry, amssymb, fancyhdr, wrapfig, epstopdf, amsmath, caption, hyperref, amsthm, bbm; the wrapper is the lane's pinned \texttt{amsart} editorial wrapper at 10pt on A4, which restates the theorem-like environments the retained text needs and adds the article's own macro definitions that the retained text needs (none), with extra wrapper packages (bm, bbm, dsfont, xspace, cancel, mathtools). The delivered numbering is the unit's own. Assets: no retained span contains a figure environment, a graphics-inclusion command, a PDF-page inclusion, a table or a TikZ picture; no image file is copied into this package, the package declares no asset dependency and no image file is redistributed with it. Licence: version arXiv:2603.01312v1 of this article is distributed under the Creative Commons Attribution 4.0 International licence, as recorded in the arXiv licence metadata for this exact version and shown on the abstract page https://arxiv.org/abs/2603.01312v1. A full rights notice is delivered at licenses/arxiv-2603.01312-CC-BY-4.0.txt. Characters: every retained excerpt is pure ASCII, so the delivered engine, XeLaTeX with its default Latin Modern text font, reports no missing character.
\begin{theoremN}
    For $\mathcal{P}_{u/v}^j(q,a,t)$ the Poincar\'e polynomial for the $j$-colored HOMFLY complex for $\tau_{u/v}$, we have
    \begin{equation}
        \mathcal{P}_{u/v}^j(q,a,t)=\mathcal{F}_{u/v}\left(\left(\widetilde{\mathcal{P}_{u/v}^1}(q,a,t)\right)^j\right).
    \end{equation}
    for all $j\geq 0$.   
\end{theoremN}

\begin{proof}
    We know that 
    \begin{multline}
        \mathcal{P}_{u/v}^j(q,a,t)=\sum_{\substack{\textbf{d}=(d_1,...,d_{u+v})\in \mathbb{N}^{u+v}\\d_1+...+d_{u+v}=j}}q^{S\cdot\bf{d} }a^{A\cdot \bf{d}}q^{\bf{d}\cdot Q\cdot \bf{d}^t} t^{T\cdot\bf{d}} {d_1+...+d_p\brack d_1,...,d_p}{d_{p+1}+...+d_{u+v}\brack d_{p+1},...,d_{u+v}}\\ \times X[d_1+...+d_{u+v},d_1+...+d_p],
    \end{multline}
    so we clearly have
    \begin{align*}
    \mathcal{P}_{u/v}^1(q,a,t) & = \sum_{i=1}^p q^{SQ_i}a^{A_i}t^{T_i}X[1,1]+\sum_{i=p+1}^{u+v} q^{SQ_i}a^{A_i}t^{T_i}X[1,0] \\
    &= \mathcal{F}_{u/v}\left(\widetilde{\mathcal{P}_{u/v}^1}(q,a,t)\right).
    \end{align*}
    Now, we proceed by induction on $j$.

    Suppose we know 
    \[
        \left(\widetilde{\mathcal{P}_{u/v}^1}(q,a,t)\right)^j=\sum_{\substack{\textbf{d}'=(d_1',...,d_{u+v}')\in \mathbb{N}^{u+v}\\d_1'+...+d_{u+v}'=j}}q^{S\cdot\bf{d}' }a^{A\cdot \bf{d}'}q^{\bf{d}'\cdot Q\cdot (\bf{d}')^t} t^{T\cdot\bf{d}'} {d_1'+...+d_p'\brack d_1',...,d_p'}{d_{p+1}'+...+d_{u+v}'\brack d_{p+1}',...,d_{u+v}'}\xi_1^{d_1'}...\xi_{u+v}^{d_{u+v}'},
    \]
    so that $\mathcal{F}_{u/v}\left(\left(\widetilde{\mathcal{P}_{u/v}^1}(q,a,t)\right)^j\right)=\mathcal{P}_{u/v}^j(q,a,t)$. In the product $\left(\widetilde{\mathcal{P}_{u/v}^1}(q,a,t)\right)^j\left(\widetilde{\mathcal{P}_{u/v}^1}(q,a,t)\right)$, we know that the term involving $\xi_1^{d_1}...\xi_{u+v}^{d_{u+v}}$ (so $d_1+...+d_{u+v}=j+1)$ is the sum of the terms coming from products of the $\xi_1^{d_1}...\xi_i^{d_i-1}...\xi_{u+v}^{d_{u+v}}$ in $\left(\widetilde{\mathcal{P}_{u/v}^1}(q,a,t)\right)^j$ with the $\xi_i$ term in $\widetilde{\mathcal{P}_{u/v}^1}(q,a,t)$. Let $\textbf{d}=(d_1,...,d_i,...,d_{u+v})$ and $\textbf{d}_i=(d_1,...,d_i-1,...,d_{u+v})$. We will assume $i\leq p$ for now; the case where $i>p$ is similar.

    Now, we multiply the $\xi_1^{d_1}...\xi_i^{d_i-1}...\xi_{u+v}^{d_{u+v}}$ and $\xi_i$ terms explicitly. We get
    \begin{align*}
    &\left(q^{SQ_i}a^{A_i}t^{T_i}\xi_i\right)\left(q^{S\cdot\textbf{d}_i }a^{A\cdot \textbf{d}_i}q^{\textbf{d}_i\cdot Q\cdot \textbf{d}_i^t} t^{T\cdot\bf{d}_i} {d_1+...+d_i-1+...+d_p'\brack d_1,...,d_i-1,...,d_p'}{d_{p+1}+...+d_{u+v}\brack d_{p+1},...,d_{u+v}}\xi_1^{d_1}...\xi_i^{d_i-1} ...\xi_{u+v}^{d_{u+v}'}\right) \\
    & =q^{S\cdot \textbf{d}}a^{A\cdot \textbf{d}}t^{T\cdot \textbf{d}}q^{Q_{ii}+\textbf{d}_i\cdot Q\cdot \textbf{d}_i^t}{d_1+...+d_i-1+...+d_p'\brack d_1,...,d_i-1,...,d_p'}{d_{p+1}+...+d_{u+v}\brack d_{p+1},...,d_{u+v}}\left(\xi_i\cdot\left(\xi_1^{d_1}...\xi_i^{d_i-1} ...\xi_{u+v}^{d_{u+v}'}\right)\right).
    \end{align*}
    Now, one can easily check by the definition of multiplication in $\widetilde{\mathbb{C}\mathcal{G}_{u/v}}$ that
    \[
    \left(\xi_1^{d_1}...\xi_i^{d_i-1} ...\xi_{u+v}^{d_{u+v}'}\right)\cdot \xi_i=q^{2\sum_{l=i+1}^p d_l}q^{2\sum_{l=1}^{u+v}d_lQ_{il}-2Q_{ii}}\left(\xi_1^{d_1}...\xi_i^{d_i} ...\xi_{u+v}^{d_{u+v}'}\right)
    \]
    Additionally,
    \[
    \textbf{d}\cdot Q\cdot \textbf{d}^t = \textbf{d}_i\cdot Q\cdot \textbf{d}_i^t+ 2\sum_{l=1}^{u+v}Q_{il}d_l-Q_{ii},
    \]
    which means
    \[
    q^{Q_{ii}+\textbf{d}_i\cdot Q\cdot \textbf{d}_i^t}q^{2\sum_{l=1}^{u+v}d_lQ_{il}-2Q_{ii}}=q^{\textbf{d}\cdot Q \cdot \textbf{d}^t}.
    \]
    Thus, if we add the $u+v$ terms that we get from the multiplication that contribute to the $\xi_1^{d_1}...\xi_{u+v}^{d_{u+v}}$ term, we get

%\begin{align*}
   %& \sum_{i=1}^p \left(q^{2\sum_{l=i+1}^p d_l}{d_1+...+d_i-1+...+d_p\brack d_1,...,d_i-1,...,d_p}{d_{p+1}+...+d_{u+v}\brack d_{p+1},...,d_{u+v}} \right)q^{S\cdot \textbf{d}}a^{A\cdot \textbf{d}}t^{T\cdot \textbf{d}}q^{\textbf{d}\cdot Q\cdot \textbf{d}^t}\xi_1^{d_1}...\xi_{u+v}^{d_{u+v}} \\
  % & + \sum_{i=p+1}^{u+v} \left(q^{2\sum_{l=i+1}^q d_l}{d_1+...+d_p\brack d_1,...,d_i,...,d_p}{d_{p+1}+...+d_i-1+...+d_{u+v}\brack d_{p+1},...,d_i-1,...,d_{u+v}} \right)q^{S\cdot \textbf{d}}a^{A\cdot \textbf{d}}t^{T\cdot \textbf{d}}q^{\textbf{d}\cdot Q\cdot \textbf{d}^t}\xi_1^{d_1}...\xi_{u+v}^{d_{u+v}} \\
  % & = q^{S\cdot \textbf{d}}a^{A\cdot \textbf{d}}t^{T\cdot \textbf{d}}q^{\textbf{d}\cdot Q\cdot \textbf{d}^t}{d_1+...+d_p\brack d_1,...,d_i,...,d_p}{d_{p+1}+...+d_i-1+...+d_{u+v}\brack d_{p+1},...,d_i-1,...,d_{u+v}}\xi_1^{d_1}...\xi_{u+v}^{d_{u+v}},
  % \end{align*}

where the second line comes from applying the same argument we used before, but when $i>p$, and the third line follows from applying Pascal's identity to the first two lines. Thus, we get 
\[
        \left(\widetilde{\mathcal{P}_{u/v}^1}(q,a,t)\right)^{j+1}=\sum_{\substack{\textbf{d}=(d_1,...,d_{u+v})\in \mathbb{N}^{u+v}\\d_1+...+d_{u+v}=j+1}}q^{S\cdot\bf{d} }a^{A\cdot \bf{d}}q^{\bf{d}\cdot Q\cdot \bf{d}^t} t^{T\cdot\bf{d}} {d_1+...+d_p\brack d_1,...,d_p}{d_{p+1}+...+d_{u+v}\brack d_{p+1},...,d_{u+v}}\xi_1^{d_1}...\xi_{u+v}^{d_{u+v}},
    \]
so the theorem follows by induction, upon applying $\mathcal{F}_{u/v}$.  
\end{proof}

\end{document}
